%%%%%%%%%%%%%%%%%%%%%%%%%%%%%%%%%%%%%%%%%
% Academic Paper Response Letter.
% LaTeX Template
% Version 1.0 (20/07/2023)
%
% This template refers to the toolbox tutorial (in Chinese):
% https://liam.page/2016/07/22/using-the-tcolorbox-package-to-create-a-new-theorem-environment/
% Please refer to the tutorial for your customization.
%
% Template license: Creative Commons CC BY 4.0
%
% Author: Yuhan Xie
% Contact email: yhxie0107@gmail.com
% Any revision advice is welcome!
%
%%%%%%%%%%%%%%%%%%%%%%%%%%%%%%%%%%%%%%%%%


\documentclass{article}
%% Language and font encodings
\usepackage[english]{babel}
\usepackage{times} % make the section title in Times New Roman

\usepackage{booktabs}
\usepackage{tabu}
\usepackage[T1]{fontenc}

%% Sets page size and margins
\usepackage[a4paper,top=2.5cm,bottom=2cm,left=1.7cm,right=1.7cm,marginparwidth=1.75cm]{geometry}

%% Useful packages
\usepackage{amsmath}
\usepackage{graphicx}
\usepackage[colorinlistoftodos]{todonotes}
\usepackage[colorlinks=true, allcolors=blue]{hyperref}
\usepackage{indentfirst}
\usepackage{graphicx}
\usepackage{subfigure}
\usepackage{float}
\usepackage{threeparttable}
\usepackage{multirow}
\usepackage{url}
\usepackage{xurl} % allow long URLs to break across lines


\usepackage[most]{tcolorbox}


%% Comment box environment used throughout the letter
\newtcbtheorem[auto counter]{cmt}{Comment}{
	colbacktitle = black!60!white, colframe = black!60!white,
	colback = black!5!white,
	separator sign={},
	fonttitle=\bfseries,%fontupper=\itshape,
}{t}


\begin{document}

%%%%%%%%%%%%%%%%%%%%%%%%%%%%%%%%%%%%%%%%%%%%%%%%%%%%%%%%%%%%%%%%%%%%%%%%%%%%%%%
% Response to the Editor: PONE-D-26-28965R1 (second revision)
%
% The decision letter (../journal_feedback.md) contains no reviewer comments,
% only the Guest Editor's points, so the letter has one section:
%   Additional Editor Comments ...  6 items
% The decision's "Journal Requirements" paragraph is boilerplate about
% reviewer-recommended citations. With no reviewer comments it needs no answer.
%
% Comment boxes quote the decision letter verbatim. The state of each response
% is marked visibly in the body of the letter, never here: %-comments do not
% reach the PDF that the authors actually read.
%%%%%%%%%%%%%%%%%%%%%%%%%%%%%%%%%%%%%%%%%%%%%%%%%%%%%%%%%%%%%%%%%%%%%%%%%%%%%%%


\begin{center}
{\Large\bfseries Response to the Editor}

\vspace{0.2cm}
Manuscript PONE-D-26-28965R1

\vspace{0.1cm}
\textit{Zero-shot generalization of EEG attention-state decoding across
increasingly stringent validation strategies on open datasets}
\end{center}

\vspace{0.3cm}

\noindent
Dear Dr. Jouzdani,

\vspace{0.2cm}
\noindent
Thank you for your careful assessment and constructive comments. We have revised
the manuscript to clarify the preprocessing and validation protocols,
participant-level uncertainty, the interpretation of transfer comparisons and
protocol-defined labels, and the ethical basis of this secondary analysis. The
revisions to the Results section clarify the presentation and interpretation of
the existing findings; no new analyses were conducted, and no experiments were
rerun or metrics recalculated. Your comments are reproduced below and followed
by our point-by-point responses.

\vspace{0.4cm}
\noindent
Yours sincerely,

\vspace{0.2cm}
\noindent
Sanzhar Sovet, Beibit Abdikenov, Manzura Zholdasova, Altyngul Kamzanova,
Medet Mukushev and Almira Kustubayeva


%%%%%%%%%%%%%%%%%%%%%%%%%%%%%%%%%%%%%%%%%%%%%%%%%%%%%%%%%% Editor comments
\newpage
\setcounter{tcb@cnt@cmt}{0}
\section*{Additional Editor Comments}

\vspace{0.4cm}

%%%%%%%%%%%% Editor Comment 1: preprocessing and data leakage
\begin{cmt}{}{}

Preprocessing and data leakage: Please specify exactly which data are used to
fit ICA and all other data-dependent preprocessing steps in each validation
protocol. Clarify whether any information from held-out participants or target
datasets contributes to preprocessing, feature transformation, model selection,
or other fitted operations. Explain how the protocol supports the claim of
zero-shot generalization without target-specific adaptation.

\end{cmt}
\vspace{0.1cm}
\noindent
\underline{\textbf{Response}}
\vspace{0.2cm}

We thank the editor for raising this point. We clarified preprocessing and
zero-shot evaluation in «Methodology» and «Validation strategies»,
and aligned «Data preparation» with those revisions.

In «Validation strategies», the opening paragraph now states the general
rule for fitting ICA in the zero-shot protocols. Each of the four protocol
descriptions then specifies which recordings were used for its ICA fits. The
baseline description distinguishes the two datasets: ICA was fitted separately
for each selected session in the mental attention-state dataset, whereas in
SAM-40 it was fitted across a participant's selected trials and could span
training and test folds. This shared decomposition may favor the SAM-40 baseline
relative to the participant-disjoint protocols. The cross-subject, cross-task,
and cross-dataset descriptions specify how ICA was fitted for held-out
recordings and confirm that source and target recordings were never pooled in
one fit.

In «Methodology», we clarify that ICA cleaning was the only fitted step
that used unlabeled target recordings. Feature transformation was applied in the
same form to source and target recordings; standardization, class weights, and
decoder parameters were determined from source training data only. Target
labels were never used, and target data played no role in model selection,
hyperparameter search, early stopping, checkpoint selection, decoder
retraining, or tuning. Thus, \textit{zero-shot} refers to applying the
source-trained decoder unchanged to target recordings, without target-guided
decoder adaptation.

\vspace{0.3cm}


%%%%%%%%%%%% Editor Comment 2: uncertainty and sample size
\begin{cmt}{}{}

Uncertainty and sample size: Please discuss the limitations imposed by the five
participants in Dataset A. Clarify the unit of resampling used for confidence
intervals and ensure that the uncertainty estimates reflect the independent
sampling units. Avoid interpreting small differences between baseline and
cross-subject estimates as evidence of equivalence.

\end{cmt}
\vspace{0.1cm}
\noindent
\underline{\textbf{Response}}
\vspace{0.2cm}

Thank you for this comment. We revised «Generalization under increasingly
stringent validation strategies» and «Performance consistency across
decoders», both subsections of «Results and analysis», as well as «Discussion».
These revisions clarify the confidence intervals, the interpretation of small
differences between baseline and cross-subject validation, and the limitation
imposed by the five participants in the mental attention-state dataset.

In «Generalization under increasingly stringent validation strategies», the
third paragraph explains that the 95\% intervals use 10,000 bootstrap resamples
of target participants, with five participants in the mental attention-state
dataset and 40 in SAM-40. The models are held fixed, and windows and folds are
not resampled independently. The fifth paragraph reports the differences of
0.1 and 2.0 percentage points between baseline and cross-subject performance as
descriptive and states that they do not establish equivalence.

In «Performance consistency across decoders», the first paragraph explains
that the intervals shown across scenarios also use 10,000 resamples of target
participants. Each resample uses the same participants for all five decoders.

In «Discussion», the fourth paragraph's second limitation explains that the
mental attention-state dataset has only five participants. This limits
population generalization and the precision of participant-level bootstrap
intervals, despite the large number of EEG windows per participant.

\vspace{0.3cm}


%%%%%%%%%%%% Editor Comment 3: comparability of validation protocols
\begin{cmt}{}{}

Comparability of validation protocols: Please explain the extent to which the
validation comparisons isolate changes in participant identity, task
composition, and dataset. Ensure that transfer results are compared with
appropriate target-specific baselines and that the conclusions do not attribute
observed differences to a single factor when several factors change
simultaneously.

\end{cmt}
\vspace{0.1cm}
\noindent
\underline{\textbf{Response}}
\vspace{0.2cm}

We thank the editor for raising this issue. We revised «Validation
strategies», «Results and analysis», and «Discussion» to clarify what each
validation comparison can show and how the transfer results should be
interpreted.

In «Validation strategies», the second paragraph distinguishes the
baseline from the three zero-shot protocols and explains what changes in each
comparison. The baseline holds out recording units but may include participants
represented in training. Cross-subject validation keeps the condition
composition unchanged, making the unseen participant the main shift. Cross-task
validation changes both the participant and task composition; comparisons
involving the full SAM-40 composition also have active tasks in common.
Cross-dataset validation changes the participant population, experimental
paradigm, recording setup, and condition definitions together. The paragraph
states that each zero-shot result is compared with the baseline for its target
composition and that differences involving cross-task or cross-dataset transfer
cannot be attributed to one factor alone.

In «Results and analysis», the sentence reporting the 6.4-point mean decrease now
specifies that this summary includes both transfers between task-specific
compositions and transfers involving the full SAM-40 composition, whose active
tasks overlap with those of the task-specific compositions.

In «Discussion», we state explicitly that the observed reductions cannot
be attributed to task composition or dataset alone because several factors
changed together.

\vspace{0.3cm}


%%%%%%%%%%%% Editor Comment 4: interpretation of attention labels
\begin{cmt}{}{}

Interpretation of attention labels: Please maintain a clear distinction between
classifying the protocol-defined conditions and decoding attention as a general
psychological construct. Discuss the potential influence of task demands and
other differences between resting and active cognitive conditions.

\end{cmt}
\vspace{0.1cm}
\noindent
\underline{\textbf{Response}}
\vspace{0.2cm}

We thank the editor for highlighting this distinction. We revised
«Methodology» and «Discussion» to make clear that the analysis classifies
protocol-defined conditions and to address differences between resting and
active cognitive conditions.

In «Methodology», the second paragraph states that the high- and
low-attention labels denote protocol-defined conditions rather than
independently validated levels of attention, and that they may reflect
differences in task demands, arousal, sensory input, stress, or motor
responses between resting and active conditions. It specifies that the
analyses therefore assess whether decoders discriminate the specified
protocol conditions, not whether they decode attention as a general
psychological construct.

In «Discussion», the fourth paragraph's first limitation now states that the
reported accuracies reflect classification of the protocol-defined
conditions rather than decoding of attention as a general psychological
construct, and reiterates that task demands, stress, arousal, sensory input,
or motor responses may also contribute, using the SAM-40 contrast of active
tasks with rest while listening to music as an example.

\vspace{0.3cm}


%%%%%%%%%%%% Editor Comment 5: ethics and reporting
\begin{cmt}{}{}

Ethics and reporting: Please clarify the ethical approval, consent, and
exemption status applicable to the original datasets and the present secondary
analysis, as appropriate. Ensure that the ethics statement is consistent with
the submission metadata.

\end{cmt}
\vspace{0.1cm}
\noindent
\underline{\textbf{Response}}
\vspace{0.2cm}

We thank the editor for this comment. In the text we added a new
«Ethics statement» subsection under «Open-access datasets»,
extended the description of each dataset to report the consent and ethical
oversight documented in the original studies, and revised the opening of the
«Methodology» section to state explicitly that this study is a
secondary analysis of two publicly available datasets.

For the mental attention-state dataset, the original study reported that
participants were healthy university student volunteers who gave written
informed consent after being informed of the study objectives and
procedures, and that the study was conducted in accordance with the ethical
guidelines of Mersin University. For SAM-40, the original study reported
informed consent from all 40 participants and approval from the Institutional
Ethics Committee of Gauhati University (GUIEC/2019/019, 15 October 2019),
with voluntary and anonymous participation in compliance with the
Declaration of Helsinki. Our «Ethics statement» makes clear that this work is a
secondary analysis of these public recordings with no new data collection or
participant contact, so no additional ethical approval or consent was
required.

\textbf{[TODO --- supervisor to confirm.]} Consistency between this statement
and the submission metadata in Editorial Manager falls under the
corresponding author's part of the submission, so we leave that confirmation
to him.

\vspace{0.3cm}


%%%%%%%%%%%% Editor Comment 6: final consistency check
\begin{cmt}{}{}

Final consistency check: Please verify that the final clean manuscript, figures,
supplementary tables, and response letter contain consistent terminology,
numerical results, and methodological descriptions.

\end{cmt}
\vspace{0.1cm}
\noindent
\underline{\textbf{Response}}
\vspace{0.2cm}

Thank you for requesting a final consistency check. We reviewed the clean
manuscript, figures, supplementary tables, and response letter. We aligned the descriptions of
sampling, ICA fitting, SAM-40 recording units, model inputs, and the aggregation
and comparison of transfer results, and checked terminology and numerical
summaries across these materials.

\vspace{0.3cm}


\end{document}
